\chapter{القول بعلم وذم التقليد وعربية القرآن}

\section{وجوب القول بعلم وذم التقليد الأعمى}
ابتدأ الإمام الشافعي رحمه الله بتأصيل عظيم، حيث قال: «فالواجب على العالِمين ألا يقولوا إلا من حيث علموا. وقد تكلم في العلم من لو أمسك عن بعض ما تكلم فيه منه، لكان الإمساك أولى به وأقرب من السلامة له إن شاء الله». فلا يجوز للعالم أن يتكلم إلا بعلم راسخ، ولا يجوز لأحد أن يُحلل أو يُحرم إلا بدليل شرعي.
ثم حذر رحمه الله من التقليد الأعمى قائلاً: «وبالتقليد أغفل من أغفل منهم، والله يغفر لنا ولهم». فلا ينبغي لطالب العلم أن يقلد تقليداً أعمى، بل يجب أن يعرف الدليل وفقه السلف لهذا الدليل؛ فإن التقليد بغير دليل شرعي ضياع، وكما قيل: «لا يقلد إلا عصبي أو غبي».

\section{مكانة لغة العرب وسعتها}
القرآن الكريم نزل بلسان عربي مبين، والعالم الذي يتصدر للفتوى والاجتهاد لا بد أن يكون عالماً بلسان العرب؛ ليعلم متى يكون اللفظ عاماً أو خاصاً، ومطلقاً أو مقيداً. والعلم بلسان العرب يستوي مع العلم بالسنن عند أهل الفقه؛ فكما أنه لا يوجد فرد جمع السنن كلها، كذلك لا يوجد من أحاط بلغة العرب كاملة إلا نبي، ولكنها لا تغيب عن مجموع الأمة. وإذا جُمع علم عامة أهل العلم، لم يفتهم شيء من السنن ولا من لغة العرب.

\section{عربية القرآن الخالصة والرد على الشبهات}
أصّل الشافعي رحمه الله لحقيقة هامة، وهي أنه لا يوجد في القرآن أي حرف غير عربي، وأن القرآن كله عربي خالص، مستدلاً بقوله تعالى: ﴿وَمَا أَرْسَلْنَا مِن رَّسُولٍ إِلَّا بِلِسَانِ قَوْمِهِ﴾. 
ورد على من زعم وجود كلمات أعجمية (مثل: استبرق، وسندس) بأن العرب قد تكلمت بها وعرّبتها قبل نزول القرآن فصارت من لغتها. والاعتزاز بالعربية هو جزء من الاعتزاز بالدين والشخصية الإسلامية، ولا يجوز لأهل لسان النبي صلى الله عليه وسلم أن يكونوا أتباعاً لأهل لسان غيره، بل يجب أن يكون لسانهم هو المتبوع.

\section{النصيحة للمسلمين ووجوه البيان في لغة العرب}
بيّن الشافعي أن توضيح لغة العرب للناس هو من باب النصيحة، والنصيحة للمسلمين فرض ودين، كما قال النبي صلى الله عليه وسلم: «الدين النصيحة».
وقد خاطب الله العرب بكتابه على ما تعرف من معانيها، ومن وجوه ذلك في لسان العرب:
\begin{itemize}
    \item \textbf{العام الظاهر:} الذي يراد به العموم.
    \item \textbf{العام المراد به الخصوص:} لفظ عام لكن المراد به خاص، ويُستدل على ذلك بسياق الكلام.
    \item \textbf{تعدد الأسماء للمسمى الواحد:} كالأسد (ليث، حفص، حيدرة).
    \item \textbf{تعدد المعاني للاسم الواحد:} كالعين (عين البصر، عين الماء، الجاسوس، الذهب).
\end{itemize}
وهذه الوجوه تعرفها العرب بالسليقة، ومَنْ جَهِلَ لسانهم قد يستنكرها، ولا عبرة بإنكار من جهل لغة العرب التي نزل بها القرآن.

\section{خطورة التكلف والقول بغير علم}
ختم الشافعي هذا المبحث بقاعدة منهجية عظيمة، فقال: «ومن تكلف ما جهل، وما لم تثبته معرفته، كانت موافقته للصواب -إن وافقه من حيث لا يعرفه- غير محمودة، وكان بخطئه غير معذور».
فمن تكلم في الدين بغير علم ولا التزام بمنهج السلف، فإن وافق الصواب صدفة فإنه لا يُحمد على ذلك لأنه لم يسلك الطريق الصحيح، وإن أخطأ فلا يُعذر بخطئه. أما المجتهد الذي يسير على منهج السلف، فإنه إن أصاب فله أجران، وإن أخطأ فله أجر وهو معذور.